\scanpage{249}\mspage{- \handwritten{6} -}
n$^{\mathrm{o}}$ \handwritten{293}

$\mathrm z$ soit l'ensemble des éléments $z\in\mathrm r$ tels que $zu\in\mathrm f$.
Posons $\mathrm r'=\mathrm o[t]+u\mathrm r$ ; on vérifie immédiatement que $\mathrm r'$
est un sous-anneau de $\mathrm r$. Si $y$ est un élément de $\mathrm o[t]$, une
condition nécessaire et suffisante pour que $y\,\mathrm r'\subset\mathrm o[t]$ est
que $yu\mathrm r\subset\mathrm o[t]$, donc que $yu\in\mathrm f$, ou encore que $y\in\mathrm z\cap\unsure{\mathrm e|\mathrm o}[t]={}
=\mathrm y$ ; $\mathrm y$ est donc bien le conducteur de $\mathrm r'$ par rapport à $\mathrm o[t]$.
Supposons pour un moment que $\omega(t)$ ne soit pas transcendant par
rapport à $\omega(\mathrm o)$ ; il y a alors un polynôme $\bar H\ne 0$ à coefficients
dans $\omega(\mathrm o)$ t\deleted{\unsure{a}}el que $\bar H(\omega(t))=0$. On en déduit tout de suite qu'il
y a un polynôme $H$ à coefficients dans $\mathrm o$ dont le coefficient dominant
$c$ n'est pas dans $\mathrm z$ tel que $H(t)\in\mathrm z$. On a alors $H(t)\in\mathrm y\unsure{'|\vphantom{'}}$ ;
il résulte alors du lemme \handwritten{3}, appliqué au couple $(\mathrm o,\mathrm r')$,
qu'il y a une puissance de $c$ qui appartient à $\mathrm y$, donc à $\mathrm z$ ; or
ceci est impossible puisque $\mathrm z$ est premier et $c\notin\mathrm z$.

\uline{Lemme} \handwritten{5} - La \uline{situation étant celle décrite dans l'énoncé du
lemme} \handwritten{3},- \uline{supposons de plus que} $\mathrm o$ \uline{soit noetherien}, \uline{Soient} $\mathrm q$ \uline{un
idéal premier maximal de} $\mathrm r$ \uline{et} $\mathrm m$ \uline{l'idéal premier} $\mathrm q\cap\mathrm o$ \uline{de} $\mathrm o$,
\uline{Si l'anneau,local de} $\mathrm q$ \uline{est distinct de celui de l'idéal} $\mathrm m$ \uline{de} $\mathrm o$,
\uline{il y a un idéal premier de} $\mathrm r$ \uline{contenu dans} $\mathrm q$, \uline{distinct de} $\mathrm q$
\uline{et contenant} $\mathrm m$.

Supposons d'abord que $\mathrm q$ contienne le conducteur $\mathrm f$ de $\mathrm r$
par rapport à $\mathrm o[t]$ ; nous allons voir qu'on est alors dans le
second cas. L'idéal \hspace{1.5em} contient un élément minimal $\mathrm r$ de l'ensemble
des idéaux premiers qui contiennent $\mathrm f$ ; $\mathrm z$ est donc un idéal
premier associé à $\mathrm f$. Soit $\omega$ l'homomorphisme canonique de $\mathrm r$
sur $\mathrm r/\mathrm z$ ; $\omega(t)$ est transcendant sur $\omega(\mathrm o)$ (lemme \handwritten{3}) et $\mathrm r/\mathrm z$
est entier sur $(\omega(\mathrm o))[\omega(t)]$. Le noyau de la restriction
% archive sheet number: 248
